\section{术语表（中英对照）}
\label{app:terms}

本表是五册手册共用的候选术语池中，本册实际用到的部分。排序按主题而非字母；同一概念只给一条主译名，括号内为可接受的别名。管理侧维护的总表若与本表冲突，以总表为准，并在回执中报差异。

\begin{center}
\footnotesize
\setlength{\tabcolsep}{3pt}
\begin{longtable}{p{0.16\linewidth}p{0.20\linewidth}p{0.33\linewidth}p{0.20\linewidth}}
\toprule
中文 & English & 简明定义 & 夜半对应 \\
\midrule
\endhead
音高 & pitch & 听觉上的高低，由基频决定 & 钢琴卷帘纵轴、音符 pitch \\
频率 & frequency & 每秒振动次数，单位 Hz & 不直接显示；导出/分析用 \\
振幅 & amplitude & 振动幅度，影响响度 & 力度、音量、推子 \\
音色 & timbre & 泛音比例与包络造成的“质感” & 乐器、合成器、效果链 \\
基频 & fundamental frequency & 一个音最主要的振动频率 & 音高行对应的实际频率 \\
泛音 & harmonic & 基频整数倍的成分 & 合成器/采样音色 \\
半音 & semitone & 十二平均律的最小音高单位 & 卷帘相邻一行 \\
全音 & whole tone & 两个半音 & 卷帘隔一行 \\
八度 & octave & 频率比 2:1 的音程 & 卷帘上下 12 行 \\
音名 & note name & C/D/E…等绝对名称 & 卷帘左侧音名 \\
科学音高记法 & scientific pitch notation & 音名＋八度数字，如 C4 & 卷帘 C4 标签 \\
MIDI 编号 & MIDI note number & 0–127 的整数音高编码 & C4=60，A4=69 \\
十二平均律 & 12-tone equal temperament & 八度均分十二份的律制 & 默认音高体系 \\
音分 & cent & 一个半音的 1/100 & 采样编辑/移调参数 \\
纯律 & just intonation & 用简单整数频率比的律制 & 不作为默认 \\
音程 & interval & 两个音之间的音高距离 & 两音符行距 \\
大音程 & major interval & 含大二/大三/大六/大七度等 & 和弦构造 \\
小音程 & minor interval & 比同名大音程少一个半音 & 和弦构造 \\
纯音程 & perfect interval & 纯一/纯四/纯五/纯八度 & 和弦构造 \\
增/减音程 & augmented/\allowbreak diminished interval & 比纯/大/小再扩或缩一个半音 & 进阶和声 \\
协和 & consonance & 听感稳定、易融合的音程 & 配器听感 \\
不协和 & dissonance & 听感紧张、需要解决的音程 & 张力音 \\
音阶 & scale & 按音级排列的音高集合 & 音阶工具 \\
调式 & mode & 以不同音级为主音的七声音阶变体 & 风格写作 \\
调 & key & 以某个主音为中心的音高体系 & 工程调性笔记 \\
调号 & key signature & 记谱时省略升降号的约定 & 读谱 \\
五声音阶 & pentatonic scale & 五个音的音阶（常见宫商角徵羽） & 流行/民谣写作 \\
和弦 & chord & 两个以上音同时/功能上成组 & 卷帘叠音符 \\
三和弦 & triad & 根音＋三音＋五音 & 常用和弦 \\
七和弦 & seventh chord & 三和弦＋七音 & 爵士/流行 \\
转位 & inversion & 改变最低音而不改和弦音 & 和弦标记斜杠 \\
根音 & root & 和弦命名的基准音 & 和弦标记第一个音 \\
和弦标记 & chord symbol & 如 Cmaj7、Am7、G7 & 工程/谱面 \\
进行 & progression & 和弦按顺序连接 & 段落写作 \\
主功能 & tonic function & 以主和弦为中心 & 调性分析 \\
属功能 & dominant function & 推向主功能的 V 级 & 终止式 \\
下属功能 & subdominant function & 准备属功能的 IV/ii 级 & 终止式 \\
终止式 & cadence & 乐句/段落的收束公式 & 曲式分析 \\
借用和弦 & borrowed chord & 从同主音其它调式借来的和弦 & 色彩和声 \\
转调 & modulation & 从一个调中心转到另一个 & 段落调性变化 \\
动机 & motif & 可重复的最小音乐细胞 & 旋律写作 \\
乐句 & phrase & 有起有收的一组动机 & 旋律写作 \\
问答句 & call and response & 前句提问、后句回答 & 旋律写作 \\
模进 & sequence & 同一动机换音高重复 & 旋律/编曲 \\
曲式 & form & 段落组织方式 & 编曲区结构 \\
主歌 & verse & 叙事段落 & 片段命名 \\
副歌 & chorus & 记忆点/高潮段落 & 片段命名 \\
桥段 & bridge & 对比与转折段落 & 片段命名 \\
五线谱 & staff & 五条横线的记谱系统 & 读谱 \\
谱号 & clef & 高音/低音等记谱基准 & 读谱 \\
升降记号 & accidental & \sharpsym / \flatsym / 还原号 & 读谱/音名 \\
力度记号 & dynamic marking & p、mf、f 等 & 力度泳道 \\
表情记号 & articulation & 连奏、断奏等 & 演奏法（采样器） \\
速度 & tempo & 每分钟拍数 BPM & 走带/BPM \\
拍 & beat & 节拍的基本脉冲 & 时间轴 \\
拍号 & time signature & 如 4/4、3/4 & 工程设置 \\
切分 & syncopation & 重音落在弱拍/弱位 & 节奏写作 \\
力度 & velocity & MIDI 音符的力度值 & 力度泳道 \\
数字音频工作站 & DAW & 录音/编辑/混音软件 & 夜半本体 \\
钢琴卷帘 & piano roll & 横时间、纵音高的音符编辑器 & P/F7 打开 \\
蓝调音阶 & blues scale & 小调五声加 \(\flat\)5 的六/七声音阶 & 风格写作 \\
蓝音 & blue note & 与调内稳定音冲突、作经过或张力的 \(\flat\)3/\(\flat\)5/\(\flat\)7 & 旋律/蓝调 \\
十二小节蓝调 & 12-bar blues & I7--IV7--V7 的十二小节骨架 & 段落写作 \\
特征音 & characteristic tone & 调式中与主音形成标志性音程的音级 & 调式写作 \\
调式循环 & modal vamp & 不靠 V--I、用主音底重复维持中心的短循环 & 段落写作 \\
声部连接 & voice leading & 各声部横向移动与纵向关系的规则 & 和弦编排 \\
再配和 & reharmonisation & 保留旋律、改变和弦解释的手法 & 和声改写 \\
低音转位 & inverted bass & 用三/五/七音作低音改变低音线 & 和弦标记斜杠 \\
drop-2 排列 & drop-2 voicing & 四声部中第二高音下移八度的开放排列 & 和弦编排 \\
复合拍 & compound meter & 每大拍含三个细分的拍号体系（6/8、12/8） & 时间轴 \\
赫米奥拉 & hemiola & 临时把两组三拍听成三组两拍的交叉 & 节奏写作 \\
动机发展 & motivic development & 模进、倒影、扩大、碎片化等变化手法 & 旋律写作 \\
段落功能 & section function & 主歌/副歌/桥段各自承担的张力与信息职责 & 编曲区结构 \\
\bottomrule
\end{longtable}
\end{center}
